\documentclass[10pt,a4paper]{article}

\def\npart {IA}
\def\nterm {Michaelmas}
\def\nyear {2014}
\def\nlecturer {M.\ G.\ Worster}
\def\ncourse {Differential Equations}
\def\ncoursehead {Differential Equations --- Long Problem}
\def\nisofficial {true}

\input{header}

\usepackage[
  a4paper,
  left=40mm,
  right=40mm,
  top=24mm,
  bottom=24mm,
  headheight=14pt,
  headsep=6mm,
  footskip=10mm,
  heightrounded
]{geometry}
\usepackage{lastpage}
\usepackage{needspace}

\hypersetup{
  pdfauthor={Adapted from notes by Dexter Chua},
  pdftitle={Differential Equations: From Local Change to Evolution},
  pdfsubject={An extended problem in differential equations}
}

\setlength{\parindent}{0pt}
\setlength{\parskip}{3pt}
\setlist[enumerate]{itemsep=2pt,topsep=4pt}
\frenchspacing
\flushbottom

\newcounter{examquestion}
\newcommand{\examquestion}[1]{%
  \Needspace{5\baselineskip}%
  \refstepcounter{examquestion}%
  \par\addvspace{0.8\baselineskip}%
  \noindent\textbf{\theexamquestion)}\enspace #1\par
}
\newenvironment{subquestions}
  {\begin{enumerate}[label=\alph*),leftmargin=2.1em]}
  {\end{enumerate}}
\newcommand{\nameditem}[2]{%
  \textbf{\textsf{#1}}\enspace\emph{#2.}\enspace
}
\newenvironment{properties}
  {\begin{enumerate}[label=(\roman*),leftmargin=2.4em]}
  {\end{enumerate}}
\newcommand{\indication}[1]{%
  \par\smallskip\noindent\emph{Indication.---} #1\par
}
\newcommand{\admitted}[3]{%
  \par\smallskip\noindent\emph{Admitted result.}\enspace
  \nameditem{#1}{#2}#3\par
}
\newcommand{\concept}[1]{%
  \par\addvspace{0.55\baselineskip}\noindent\textbf{#1.}\enspace
}
\newcommand{\exampart}[2]{%
  \clearpage
  \markboth{Part #1}{}
  \begin{center}
    {\large\bfseries Part #1. --- #2}
  \end{center}
}
\newcommand{\subpart}[1]{%
  \begin{center}
    \textbf{#1}
  \end{center}
}

\pagestyle{fancy}
\fancyhf{}
\setlength{\headwidth}{\textwidth}
\lhead{\emph{\nouppercase{\leftmark}}}
\rhead{\ncoursehead}
\cfoot{\thepage\ of \pageref{LastPage}}

\begin{document}
\pagenumbering{roman}

\begin{titlepage}
  \thispagestyle{empty}
  \begin{center}
    \vspace*{1.3cm}
    {\Large\bfseries DIFFERENTIAL EQUATIONS\par}
    \vspace{0.35cm}
    \rule{0.72\textwidth}{0.6pt}
    \vspace{0.65cm}

    {\huge\bfseries From Local Change to Evolution\par}
    \vspace{0.2cm}
    {\Large Calculus, Dynamics, and Waves\par}

    \vfill

    {\Large Extended mathematical problem\par}
    \vspace{0.3cm}
    {\large Part I\par}

    \vfill

    \begin{minipage}{0.72\textwidth}
      \centering
      \textbf{Extended paper --- intended for several sittings}\par
      \vspace{0.25cm}
      Calculators and electronic aids are not required.
    \end{minipage}

    \vfill

    \begin{minipage}{0.82\textwidth}
      \centering\small
      Adapted from the \emph{Differential Equations} notes based on lectures
      by M.\ G.\ Worster. Original notes by Dexter Chua.
    \end{minipage}
    \vspace*{0.8cm}
  \end{center}
\end{titlepage}

\pagenumbering{arabic}

\begin{center}
  {\large\bfseries Instructions and planned structure}
\end{center}

The completed paper will consist of one problem divided into the following
ten parts.  This instalment contains Part I.

\begin{enumerate}[label=\Roman*.,leftmargin=3.1em,itemsep=1pt]
  \item Differentiation, asymptotics, and local approximation;
  \item integration and multivariable differentiation;
  \item first-order linear differential equations;
  \item nonlinear first-order equations and discrete dynamics;
  \item homogeneous second-order equations and phase space;
  \item forced and equidimensional equations, recurrences, damping, and
  impulses;
  \item series solutions near ordinary and singular points;
  \item gradients, Taylor forms, and stationary points;
  \item differential systems and phase portraits;
  \item characteristics, waves, and diffusion.
\end{enumerate}

Question numbers will remain continuous when later parts are added.  The
result of any question may be used later, even if that question has not
been answered, provided its number is cited clearly.  Indications are
optional.  Definitions and notation are introduced when they first become
necessary.  No result may be used unless it belongs to the calculus background
specified below, has been stated as an admitted result, or has been proved in
an earlier question.

\exampart{I}{Differentiation, asymptotics, and local approximation}

As in the original notes, the reader is assumed to know elementary
single-variable real calculus: real limits and their algebra, continuity, the
usual elementary functions and their derivatives, the mean-value theorem,
and mathematical induction.  The definitions and results emphasized by the
course are nevertheless stated and developed here.

\subpart{A. --- Derivatives and equations of evolution}

\concept{Derivative and differentiability}
Let $I\subseteq\R$ be an open interval, let $f\colon I\to\R$, and let
$x_0\in I$.  The \emph{derivative} of $f$ at $x_0$ is
\[
  f'(x_0)=\lim_{h\to0}\frac{f(x_0+h)-f(x_0)}{h},
\]
provided this real limit exists.  In that case, $f$ is
\emph{differentiable at $x_0$}.  It is differentiable on $I$ if it is
differentiable at every $x_0\in I$.  The right-hand and left-hand derivatives
are defined by replacing $h\to0$ with $h\to0^+$ and $h\to0^-$,
respectively.  Thus $f$ is differentiable at $x_0$ precisely when both
one-sided derivatives exist and are equal.

\examquestion{\label{q:derivative}
\begin{subquestions}
  \item \nameditem{DQPL}{Difference Quotient for a Power Law}
  Define $p\colon\R\to\R$ by $p(x)=x^2$ for every $x\in\R$.
  Starting from the difference quotient, prove that $p'(x_0)=2x_0$
  for every $x_0\in\R$.

  \item \nameditem{ONDD}{One-sided Non-Differentiability Demonstration}
  Define $a\colon\R\to\R$ by $a(x)=|x|$ for every $x\in\R$.
  Compute the right-hand and left-hand derivatives of $a$ at $0$ and
  prove that $a$ is not differentiable there.

  \item \nameditem{HDC}{Higher Derivative Conventions}
  When $f$ is differentiable, its derivative function is also written
  \[
    \frac{\d f}{\d x}=f'(x)=\frac{\d}{\d x}f(x).
  \]
  If $f'$ is differentiable, define
  $f''=(f')'$ and write
  $\d^2f/\d x^2=f''(x)$.
  Let $f\colon I\to\R$ be differentiable, put
  $J=\{t\in\R:2t\in I\}$, and define $g\colon J\to\R$ by
  $g(t)=f(2t)$.  Use the difference quotient to prove
  $g'(t)=2f'(2t)$ for every $t\in J$.
\end{subquestions}}

\concept{Differential and difference equations}
Let $I\subseteq\R$ be an interval.  A \emph{differential equation} is an
equation imposed on an unknown function $x\colon I\to\R$ and some of its
derivatives.  A \emph{solution on $I$} is a function possessing the required
derivatives and satisfying the equation at every $t\in I$.  When the argument
is time $t$, write
\[
  \dot x(t)=\frac{\d x}{\d t}(t),
  \qquad
  \ddot x(t)=\frac{\d^2x}{\d t^2}(t).
\]
An \emph{initial condition} prescribes one or more values of the unknown
function and its derivatives at a specified time.

\examquestion{\label{q:differential-models}
Let $m\in\R$ satisfy $m>0$.  In a one-dimensional Newtonian model,
$t\in\R$ denotes time, a twice differentiable function
$x\colon\R\to\R$ gives the particle's position, and its force is the
function $F_x\colon\R\to\R$ defined by
\[
  F_x(t)=m\ddot x(t)\qquad(t\in\R).
\]
\begin{subquestions}
  \item \nameditem{PADE}{Prescribed Acceleration Differential Equation}
  Suppose that $F_x(t)=t^2-1$ for every $t\in\R$.  Find all possible
  position functions, or equivalently all twice differentiable functions
  $x\colon\R\to\R$ satisfying
  \[
    m\ddot x(t)=t^2-1\qquad(t\in\R).
  \]
  Prove that they are exactly the functions
  \[
    x(t)=\frac{t^4}{12m}-\frac{t^2}{2m}+c_1t+c_0
    \qquad(t\in\R),
  \]
  where $c_0,c_1\in\R$.

  \item \nameditem{PDNDE}{Position-Dependent Nonlinear Differential Equation}
  Suppose instead that $F_x(t)=x(t)^2-1$ for every $t\in\R$, so that
  \[
    m\ddot x(t)=x(t)^2-1\qquad(t\in\R).
  \]
  Explain why the occurrence of $x(t)^2$ makes this equation nonlinear.
  Determine all constant solutions $x\colon\R\to\R$.

  \item \nameditem{ICSU}{Initial Conditions Select Uniquely}
  Let $x_0,v_0\in\R$.  Among the solutions in part (a), find the unique
  one satisfying $x(0)=x_0$ and $\dot x(0)=v_0$.
\end{subquestions}}

\indication{For the converse statement in part (a), subtract the displayed
candidate from an arbitrary solution and use the standard consequence of the
mean-value theorem that a function with zero second derivative is affine.}

Write $\N_0=\{0,1,2,\ldots\}$.  A real sequence is a map
$F\colon\N_0\to\R$, whose value at $n\in\N_0$ is denoted by $F_n$.
A \emph{difference equation} is a relation between terms of an unknown
sequence.  A solution is a sequence satisfying that relation for every index
for which it is stated.

\examquestion{\label{q:difference-equations}
Consider the Fibonacci difference equation
\[
  F_{n+2}-F_{n+1}-F_n=0\qquad(n\in\N_0).
\]
\begin{subquestions}
  \item \nameditem{FRCS}{Fibonacci Recurrence Computes Successors}
  For the solution with $F_0=0$ and $F_1=1$, compute
  $F_2,F_3,\ldots,F_7$.

  \item \nameditem{TRIVFS}{Two Recurrence Initial Values Fix the Sequence}
  Let $F,G\colon\N_0\to\R$ be two solutions.  Prove by induction that
  $F_0=G_0$ and $F_1=G_1$ imply $F_n=G_n$ for every $n\in\N_0$.

  \item \nameditem{RANC}{Recurrence Alone is Not Complete}
  Exhibit two distinct solutions of the difference equation.  Explain why the
  recurrence without two initial values does not specify a unique sequence.
\end{subquestions}}

The explicit formula for the Fibonacci solution will be derived in Part VI.
These first examples also illustrate a recurring method: one may guess a
candidate solution, but the candidate becomes a solution only after its
domain, differentiability, and governing equation have been checked.

\subpart{B. --- Small and bounded asymptotic orders}

\concept{Asymptotic notation}
Let $I\subseteq\R$ be an interval, let
$x_0\in\R\cup\{-\infty,+\infty\}$ be a limit point of $I$ in the
relevant direction, and let $f,g\colon I\to\R$.  All limits below are taken
through points $x\in I$.  A \emph{deleted neighbourhood} of a finite $x_0$
consists of the points satisfying $0<|x-x_0|<\delta$ for some
$\delta>0$; at $+\infty$ it consists of the points satisfying $x>M$ for
some $M\in\R$, and at $-\infty$ it is defined analogously.

Suppose first that $g(x)\ne0$ throughout some deleted neighbourhood of
$x_0$.  We write
\[
  f(x)=o(g(x))\quad\hbox{as }x\to x_0
\]
if
\[
  \lim_{x\to x_0}\frac{f(x)}{g(x)}=0.
\]
We write
\[
  f(x)=O(g(x))\quad\hbox{as }x\to x_0
\]
if there exist $C\in\R$ with $C\geq0$ and a deleted neighbourhood of
$x_0$ throughout which
\[
  |f(x)|\leq C|g(x)|.
\]
The latter definition remains meaningful without the non-vanishing
assumption on $g$.  Informally, $o(g)$ is negligible compared with $g$,
whereas $O(g)$ is bounded in magnitude by a constant multiple of $g$.

\examquestion{\label{q:asymptotic-examples}
\begin{subquestions}
  \item \nameditem{SROZ}{Square-Root Order near Zero}
  On the interval $(0,\infty)$, prove that
  $x=o(\sqrt{x})$ as $x\to0^+$.

  \item \nameditem{SRGI}{Square Root Growth at Infinity}
  On the interval $(0,\infty)$, prove that
  $\sqrt{x}=o(x)$ as $x\to+\infty$.

  \item \nameditem{LSOB}{Local Sine Order Bound}
  Using the elementary inequality $|\sin u|\leq|u|$ for every $u\in\R$,
  prove that
  $\sin(2x)=O(x)$ as $x\to0$.

  \item \nameditem{SBI}{Sine Bounded at Infinity}
  Prove that $\sin(2x)=O(1)$ as $x\to+\infty$, although
  $\sin(2x)$ has no limit as $x\to+\infty$.

  \item \nameditem{VRGBR}{Vanishing Ratio Gives a Bounded Ratio}
  Let $f,g\colon I\to\R$ satisfy the hypotheses in the definition of
  small order.  Prove that
  \[
    f(x)=o(g(x))\quad\Longrightarrow\quad f(x)=O(g(x))
    \qquad(x\to x_0).
  \]

  \item \nameditem{ONCRF}{Order Notation denotes a Class Rather than a Function}
  Explain why $o(g)$ denotes a class of functions rather than one function.
  Interpret the conventional equality $f=o(g)$ as a membership assertion,
  and state that assertion using the symbol $\in$.
\end{subquestions}}

\examquestion{\label{q:asymptotic-algebra}
In every part, all functions are real-valued and all asymptotic statements
refer to the same limit $x\to x_0$.
\begin{subquestions}
  \item \nameditem{SRAC}{Scaling Rules for Asymptotic Classes}
  Let $\lambda\in\R$.  Prove that $f=o(g)$ implies
  $\lambda f=o(g)$, and that $f=O(g)$ implies $\lambda f=O(g)$.

  \item \nameditem{SAR}{Sum of Asymptotic Remainders}
  Prove that $f_1=o(g)$ and $f_2=o(g)$ imply
  $f_1+f_2=o(g)$.  Prove the analogous assertion with $O$ in place
  of $o$.

  \item \nameditem{PROB}{Product Remainder Order Bound}
  Prove that $f=O(g)$ and $p=O(q)$ imply $fp=O(gq)$.
  Deduce that, as $h\to0$, the relations
  $r(h)=o(h)$ and $s(h)=O(h)$ imply
  $r(h)s(h)=o(h^2)$.

  \item \nameditem{CRCL}{Composition Remainder Control Lemma}
  Let $\delta>0$, let $k\colon(-\delta,\delta)\to\R$ satisfy
  $k(h)\to0$ and $k(h)=O(h)$ as $h\to0$, and let
  $\varepsilon$ be a real-valued function defined on an open interval
  containing $0$, with $\varepsilon(u)\to0$ as $u\to0$.
  Prove
  \[
    k(h)\varepsilon(k(h))=o(h)\qquad(h\to0).
  \]
\end{subquestions}}

\examquestion{\label{q:linear-approximation}
Let $I\subseteq\R$ be an open interval, let $x_0\in I$, and let
$f\colon I\to\R$.
\begin{subquestions}
  \item \nameditem{DILA}{Derivative Implies Linear Approximation}
  Suppose that $f$ is differentiable at $x_0$.  Starting from the
  derivative limit, prove
  \[
    f(x_0+h)=f(x_0)+f'(x_0)h+o(h)
    \qquad(h\to0).
  \]

  \item \nameditem{LADD}{Linear Approximation Determines the Derivative}
  Conversely, suppose that $A\in\R$ and a function $r$, defined near
  $0$, satisfy
  \[
    f(x_0+h)=f(x_0)+Ah+r(h),
    \qquad r(h)=o(h)\quad(h\to0).
  \]
  Prove that $f$ is differentiable at $x_0$ and $f'(x_0)=A$.

  \item \nameditem{LAC}{Linear Approximation Coefficient}
  Prove that at most one number $A\in\R$ can occur as the coefficient of
  $h$ in an expansion of the form in part (b).

  \item \nameditem{QEO}{Quadratic Error Order}
  For $f(x)=x^2$ on $\R$, show directly that
  \[
    f(x_0+h)=f(x_0)+2x_0h+h^2.
  \]
  Prove that $h^2=O(h^2)$ and $h^2=o(h)$ as $h\to0$.
\end{subquestions}}

\subpart{C. --- Product, composition, and repeated differentiation}

\examquestion{\label{q:product-rule}
Let $I\subseteq\R$ be an open interval, let $x_0\in I$, and let
$u,v\colon I\to\R$ be differentiable at $x_0$.
\begin{subquestions}
  \item \nameditem{PRE}{Product Remainder Expansion}
  Introduce functions $r_u,r_v$, defined near $0$, such that
  $r_u(h)=o(h)$, $r_v(h)=o(h)$, and
  \[
    u(x_0+h)=u(x_0)+hu'(x_0)+r_u(h),
  \]
  \[
    v(x_0+h)=v(x_0)+hv'(x_0)+r_v(h).
  \]
  Multiply these two identities.

  \item \nameditem{CTRC}{Cross-Term Remainder Control}
  Using Question \ref{q:asymptotic-algebra}, prove that all terms in the
  product from part (a), other than the constant term and the terms linear in
  $h$, have sum $o(h)$.

  \item \nameditem{DPR}{Derivative Product Rule}
  Apply Question \ref{q:linear-approximation} to prove that
  \[
    (uv)'(x_0)=u'(x_0)v(x_0)+u(x_0)v'(x_0).
  \]
\end{subquestions}}

\examquestion{\label{q:chain-rule}
Let $I,J\subseteq\R$ be open intervals, let $g\colon I\to J$, let
$F\colon J\to\R$, and let $x_0\in I$.  Suppose that $g$ is
differentiable at $x_0$ and $F$ is differentiable at
$y_0=g(x_0)$.
\begin{subquestions}
  \item \nameditem{ICI}{Inner Composition Increment}
  For all sufficiently small $h\in\R$, put
  $k(h)=g(x_0+h)-g(x_0)$.  Prove that
  \[
    k(h)=hg'(x_0)+o(h),\qquad k(h)=O(h),\qquad k(h)\to0
    \quad(h\to0).
  \]

  \item \nameditem{ORF}{Outer Remainder Function}
  For every sufficiently small $s\in\R$ with $y_0+s\in J$, define
  \[
    \varepsilon(s)=
    \begin{cases}
      \dfrac{F(y_0+s)-F(y_0)-F'(y_0)s}{s},&s\ne0,\\[6pt]
      0,&s=0.
    \end{cases}
  \]
  Prove that $\varepsilon(s)\to0$ as $s\to0$.

  \item \nameditem{CRB}{Composed Remainder Bound}
  Prove
  \[
    F(g(x_0+h))-F(g(x_0))
      =F'(y_0)k(h)+k(h)\varepsilon(k(h)).
  \]
  Use Question \ref{q:asymptotic-algebra}(d) to obtain a linear
  approximation in $h$.

  \item \nameditem{DCR}{Derivative Chain Rule}
  Conclude that $F\circ g$ is differentiable at $x_0$ and
  \[
    (F\circ g)'(x_0)=F'(g(x_0))g'(x_0).
  \]
\end{subquestions}}

\concept{Repeated derivatives and binomial coefficients}
For $n\in\N_0$, set $f^{(0)}=f$ and, whenever the required derivatives
exist, define $f^{(n+1)}=(f^{(n)})'$.  Define $0!=1$ and
$n!=1\cdot2\cdots n$ for $n\in\N$.  If $n\in\N_0$ and
$r\in\N_0$ satisfies $0\leq r\leq n$, define
\[
  \binom nr=\frac{n!}{r!(n-r)!}.
\]

\examquestion{\label{q:leibniz}
Let $I\subseteq\R$ be an open interval.  The functions below are assumed to
possess every derivative that occurs.
\begin{subquestions}
  \item \nameditem{ZOPC}{Zeroth and Onefold Product Cases}
  For $u,v\colon I\to\R$, verify that
  \[
    (uv)^{(n)}
      =\sum_{r=0}^{n}\binom nr u^{(r)}v^{(n-r)}
      \tag{$L_n$}
  \]
  holds for $n=0$ and $n=1$.

  \item \nameditem{DIH}{Differentiated Induction Hypothesis}
  Let $n\in\N$.  Assume $(L_n)$ and differentiate its two sides.
  Use the product rule to obtain two finite sums representing
  $(uv)^{(n+1)}$.

  \item \nameditem{PCR}{Pascal Combination after Reindexing}
  Reindex one of the sums from part (b), and prove directly from the
  factorial definition that
  \[
    \binom nr+\binom n{r-1}=\binom{n+1}r
    \qquad(1\leq r\leq n).
  \]
  Combine the two sums, treating the terms with $r=0$ and $r=n+1$
  separately.

  \item \nameditem{GLDFP}{General Leibniz Derivative Formula for Products}
  Use induction to prove that, for every $n\in\N_0$ and all
  $n$-times differentiable functions $u,v\colon I\to\R$,
  \[
    (uv)^{(n)}(x)
      =\sum_{r=0}^{n}\binom nr
        u^{(r)}(x)v^{(n-r)}(x)
    \qquad(x\in I).
  \]
\end{subquestions}}
\indication{In part (c), replace $r+1$ by a new index in the sum containing
$u^{(r+1)}v^{(n-r)}$ before combining coefficients.}

\subpart{D. --- Taylor polynomials and formal series}

\concept{Taylor polynomial}
Let $n\in\N_0$, let $I\subseteq\R$ be open, let $x_0\in I$, and let
$f\colon I\to\R$ possess derivatives through order $n$ at $x_0$.
The \emph{Taylor polynomial of order $n$ about $x_0$} is
\[
  T_{n,x_0}f(x)
    =\sum_{r=0}^{n}\frac{f^{(r)}(x_0)}{r!}(x-x_0)^r
    \qquad(x\in\R).
\]

\admitted{TTPR}{Taylor Theorem with Peano Remainder}{Let $n\in\N_0$, let
$I\subseteq\R$ be an open interval, let $x_0\in I$, and let
$f\colon I\to\R$.  Suppose that the derivatives
$f^{(0)},f^{(1)},\ldots,f^{(n)}$ exist on a neighbourhood of $x_0$ and
that $f^{(n)}$ is continuous at $x_0$.  Then, as $h\to0$ with
$x_0+h\in I$,
\[
  f(x_0+h)
    =\sum_{r=0}^{n}\frac{f^{(r)}(x_0)}{r!}h^r+o(h^n).
\]
If, in addition, $f^{(n+1)}$ exists and is bounded on a neighbourhood of
$x_0$, the remainder after the displayed finite sum is $O(h^{n+1})$.}

\examquestion{\label{q:taylor}
Retain the notation and hypotheses of the admitted Taylor theorem.
\begin{subquestions}
  \item \nameditem{FTPC}{Finite Taylor Polynomial Construction}
  Prove that, for every integer $k\in\N_0$ with $0\leq k\leq n$,
  \[
    \bigl(T_{n,x_0}f\bigr)^{(k)}(x_0)=f^{(k)}(x_0).
  \]

  \item \nameditem{EVTF}{Equivalent Variable Taylor Form}
  Put $x=x_0+h$ and rewrite the admitted result as an expansion of
  $f(x)$ as $x\to x_0$.  State the remainder in both the $o$ form and,
  under the additional hypothesis, the $O$ form.

  \item \nameditem{LNAL}{Local Nature and Analytic Limitation}
  Explain precisely why a statement as $x\to x_0$ controls $f$ only near
  $x_0$.  Explain also why the admitted theorem, for one fixed finite
  $n\in\N_0$, does not assert that an infinite series converges to $f(x)$.

  \item \nameditem{CTE}{Cubic Taylor Expansion}
  Define $f\colon\R\to\R$ by $f(x)=x^3$.  Compute
  $T_{3,1}f(x)$ and verify directly that it equals $f(x)$ for every
  $x\in\R$.
\end{subquestions}}

\concept{Taylor series}
Let $I\subseteq\R$ be open, let $x_0\in I$, and let
$f\colon I\to\R$ possess derivatives of every order at $x_0$.
The \emph{Taylor series of $f$ about $x_0$} is the formal power series
\[
  \sum_{n=0}^{\infty}
    \frac{f^{(n)}(x_0)}{n!}(x-x_0)^n.
\]
The Taylor series about $x_0=0$ is also called the \emph{Maclaurin
series}.  The word \emph{formal} means that this definition alone asserts
neither convergence of the series nor equality of its sum with $f(x)$.

\examquestion{\label{q:taylor-series}
\begin{subquestions}
  \item \nameditem{PTSF}{Polynomial Taylor Series is Finite}
  Let $m\in\N_0$, let $a_0,\ldots,a_m\in\R$, and define
  $p\colon\R\to\R$ by
  \[
    p(x)=\sum_{r=0}^{m}a_rx^r.
  \]
  Prove that the Taylor series of $p$ about $0$ has no non-zero terms
  after degree $m$ and is equal to $p(x)$ for every $x\in\R$.

  \item \nameditem{FSD}{Formal Series Distinction}
  State separately the following three assertions about a Taylor series:
  its coefficients are defined; the resulting numerical series converges at
  a specified $x\in I$; its sum at that $x$ equals $f(x)$.  Explain which
  assertion follows from the definition alone.

  \item \nameditem{MTRF}{Maclaurin Terms for a Rational Function}
  Define $f\colon\R\setminus\{1\}\to\R$ by
  $f(x)=(1-x)^{-1}$.  Prove that $f^{(n)}(0)=n!$ for every
  $n\in\N_0$ and hence determine its formal Maclaurin series.
\end{subquestions}}

\subpart{E. --- Quotients with vanishing numerator and denominator}

An expression $f(x)/g(x)$ has the \emph{indeterminate form $0/0$} at
$x_0\in\R$ when $f(x)\to0$ and $g(x)\to0$ as $x\to x_0$.
The phrase does not assert that the quotient has a limit.

\examquestion{\label{q:local-lhopital}
Let $I\subseteq\R$ be an open interval, let $x_0\in I$, and let
$f,g\colon I\to\R$ be differentiable at $x_0$.  Suppose
$f(x_0)=g(x_0)=0$ and $g'(x_0)\ne0$.
\begin{subquestions}
  \item \nameditem{FLQE}{First-order Linear Quotient Expansions}
  Use Question \ref{q:linear-approximation} to write expansions of
  $f(x_0+h)$ and $g(x_0+h)$ to first order, with remainders in
  $o(h)$.

  \item \nameditem{PNNZD}{Punctured-Neighbourhood Non-Zero Denominator}
  Prove that there exists $\delta>0$ such that
  $g(x_0+h)\ne0$ whenever $0<|h|<\delta$.

  \item \nameditem{LQDL}{Local Quotient Derivative Limit}
  Divide the expansions from part (a) by $h$, justify the quotient using
  part (b), and prove
  \[
    \lim_{x\to x_0}\frac{f(x)}{g(x)}
      =\frac{f'(x_0)}{g'(x_0)}.
  \]
\end{subquestions}}
\indication{For part (b), use
$g(x_0+h)/h\to g'(x_0)\ne0$ and the fact that a function with a
non-zero limit is non-zero sufficiently near the limit point.}

\admitted{GDQLR}{General Derivative Quotient Limit Rule}{Let
$x_0\in\R$, let $\delta>0$, put
$I=(x_0-\delta,x_0+\delta)$, and let $f,g\colon I\to\R$ be continuous.
Suppose that $f$ and $g$ are differentiable at every
$x\in I\setminus\{x_0\}$, that
\[
  f(x_0)=g(x_0)=0,\qquad
  g'(x)\ne0\quad(x\in I\setminus\{x_0\}),
\]
and that there exists $L\in\R$ such that
\[
  \lim_{x\to x_0}\frac{f'(x)}{g'(x)}=L.
\]
Then $g(x)\ne0$ throughout some deleted neighbourhood of $x_0$ and
\[
  \lim_{x\to x_0}\frac{f(x)}{g(x)}=L.
\]}

\examquestion{\label{q:lhopital-applications}
\begin{subquestions}
  \item \nameditem{SQR}{Sine Quotient Recovery}
  Apply Question \ref{q:local-lhopital} to
  $f(x)=\sin x$ and $g(x)=x$ on $\R$ to recover
  \[
    \lim_{x\to0}\frac{\sin x}{x}=1.
  \]

  \item \nameditem{CSOQL}{Cosine Second-Order Quotient Limit}
  Explain why Question \ref{q:local-lhopital} does not apply directly to
  $(1-\cos x)/x^2$ at $0$.  Use the admitted general rule and part (a)
  to prove
  \[
    \lim_{x\to0}\frac{1-\cos x}{x^2}=\frac12.
  \]

  \item \nameditem{RDQE}{Repeated Derivative Quotient Evaluation}
  Use the admitted general rule and part (b) to prove
  \[
    \lim_{x\to0}\frac{x-\sin x}{x^3}=\frac16.
  \]

  \item \nameditem{HRVA}{Hypothesis Review before Valid Application}
  Define $f,g\colon\R\to\R$ by $g(x)=x$ for every $x\in\R$,
  $f(0)=0$, and
  \[
    f(x)=x^2\sin(1/x)\qquad(x\in\R\setminus\{0\}).
  \]
  Prove directly that $f(x)/g(x)\to0$ as $x\to0$.  Compute
  $f'(x)/g'(x)$ for $x\ne0$ and show that this derivative quotient has
  no limit as $x\to0$.  Identify the failed hypothesis of the admitted
  general rule and explain why the original quotient may nevertheless
  converge.
\end{subquestions}}

\hfill\emph{End of Part I}

\end{document}
